\lecture{21}{Thu 17 Apr 2025 14:02}{Hydrogen Atom Work}
Recall we are working with the hydrogen atom, with an electron $r_e$ and a proton $r_p$. The Coulomb potential is translation invariant
\[
	V=- \frac{e^2}{4\pi \varepsilon _0\left| r_e-r_p \right| } 
.\]
That is, moving the \emph{entire wavefunction} somewhere else preserves physics. The Schrodinger equation is given in the sum of the energies
\[
	H=\frac{P^2_p}{2m_p} +\frac{P^2_e}{2m_e} +V(|R_e-R_p|)
.\]
We can do a change of basis to write $\psi (r_e,r_p)$ in 3 coordinates. We write $\psi $ in terms of the relative position $R_{\text{rel} } =R_e-R_p$. However, since $R_{\text{rel} } $ only has 3 variables, we need 3 more for $\psi $. We use the center of mass $R_{\text{cm} } $:
\[
	R_{\text{cm} } =\frac{m_eR_e+m_pR_p}{m_e+m_p} 
.\]
Remark: $\ket{r_e,r_p} $ is indeed a tensor product $\ket{r_e} \otimes \ket{r_p} $, which is necessary since we are working in 2 particles. So $P_p$ is like $1\otimes P$, and vice versa for $P_e$.

We should be able to write $H$ as $H_{\text{cm} } +H_{\text{rel} } $, since we should be able to write it as a hamiltonian for the entire atom (center of mass), and a hamiltonian for the internal structure (relative position).

We have $P_{\text{cm} } =P_e+P_p$, which can be verified easily by showing $\left[ X_{\text{cm} } ,P_{\text{cm} }  \right] =i\hbar $ in this definition, where
\[
	P_{x,\text{cm} } =-i\hbar \left( \frac{\partial }{\partial x_e} +\frac{\partial }{\partial x_p}  \right) 
.\]

The relative momentum is given by
\[
	P_{\text{rel} } =\frac{m_eP_e-m_pP_p}{m_e+m_p} 
.\]
This gives a change of variables $\nabla_e,\nabla _p \to \nabla_{\text{cm} } ,\nabla_{\text{rel} } $. This definition of $P_{\text{rel} } $ ensures this satisfies the chain rule; in that $P_{\text{rel} } =-i\hbar \nabla_{\text{rel} } $.

Returning to our original hamiltonian
\[
	H=\frac{P^2_p}{2m_p} +\frac{P^2_e}{2m_e} +V\left( R_{\text{rel} }  \right) ,
\]
this turns into
\[
	H=\frac{P^2_{\text{cm} } }{2M} +\frac{P^2_{\text{rel} } }{2\mu } +V\left( R_{\text{rel} }  \right) ,
\]
where the reduced mass $\mu $ is
\[
	\frac{1}{\frac{1}{m_e} +\frac{1}{m_p} } 
.\]
This can be brute forced by plugging in the substitution. The hamiltonian now becomes two parts, the center of mass hamiltonian $P^2_{\text{cm} } /2M$, and the relative motion part which is the rest. The hamiltonians commute, so they can be diagonalized independently. WLOG, we can thus write $\psi (r_{\text{rel} } ,r_{\text{cm} } )$, which can be further written as $\psi_{\text{rel} } (r_{\text{rel} } )\otimes \psi _{\text{cm} } (r_{\text{cm} } )$. The tensor product then collapses to
\[
	H\ket{\psi } =\left( 1\otimes H_{\text{cm} }+H_{\text{rel} } \otimes 1  \right) )\left( \ket{\psi_{\text{cm} }  }\otimes \ket{\psi _{\text{rel} } }  \right) =\left( E_{\text{cm} } +E_{\text{rel} }  \right) \left( \ket{\psi _{\text{cm} } } \otimes \ket{\psi _{\text{rel} } }  \right) 
.\]
Solving the first equation
\[
	H_{\text{cm} } \ket{\psi _{\text{cm} } } =E_{\text{cm} } \ket{\psi _{\text{cm} } } ,
\]
we write
\[
\left( \frac{-\hbar ^2}{2m} \nabla ^2_{\text{cm} }  \right) \psi_{\text{cm}} (r_{\text{cm} }  )=E_{\text{cm} } \psi _{\text{cm} } (r_{\text{cm} } )
.\]
If we assume separability, this is solved by exponentials. This can be shown to be solved by
\[
	\frac{1}{\left( \sqrt{2\pi \hbar }  \right) ^3} \exp \left( \frac{i}{\hbar } \left( p_xx+p_yy+p_zz \right)  \right) =\frac{1}{\left( \sqrt{2\pi \hbar }  \right) ^3} \exp\left( \frac{i \mathbf{p} \cdot \mathbf{r} }{\hbar } \right)
,\]
where the solutions to the separable equations are
\[
	\frac{1}{\sqrt{2\pi \hbar } } e^{ip_{x_i} x_i} 
.\]

We've solved $\ket{\psi_{\text{cm} } } $, and now we have to real with the relative motion. Let
\[
	H_{\text{rel} } \ket{\psi _{\text{rel} } } =E_{\text{rel} } \ket{\psi _{\text{rel} } } 
.\]
In position space, this looks like
\[
	\left(\frac{\hbar ^2}{2\mu } \nabla_{\text{rel} }^2+V(R_{\text{rel}}) \right)\psi _{\text{rel} }(r_{\text{rel}} ) =E_{\text{rel} } \psi _{\text{rel} }(r_{rel} )
.\]
We will drop the ``rel'' subscripts, since this is the final equation we will solve this semester.

To solve this, we need to use spherical harmonics, since there is also going to be spherical symmetry. We need to take $(x,y,z)$ and map it to $(r,\theta ,\phi )$. The change of variables is
\[
	x=r \sin \theta \cos \phi ,
\]
\[
	y=r \sin \theta \sin \phi ,
\]
\[
	z=r \cos \theta 
.\]
This will allow us to separate it into a solution depending on $r$, and a solution depending on $(\theta ,\phi )$, allowing it to be solved.

First: what is $\nabla ^2$ in spherical coordinates? We have
\[
	\nabla ^2=\frac{1}{r^2} \left( \frac{\partial }{\partial r} \left( r^2 \frac{\partial }{\partial r}  \right) +\frac{1}{\sin \theta } \frac{\partial }{\partial \theta } \left( \sin \theta \frac{\partial }{\partial \theta }  \right) +\frac{1}{\sin ^2\theta } \frac{\partial ^2}{\partial \phi ^2}  \right) 
.\]
Let's define
\[
	\hat{\nabla}^2 \coloneqq \frac{1}{\sin \theta } \frac{\partial }{\partial \theta } \left( \sin \theta \frac{\partial }{\partial \theta }  \right) +\frac{1}{\sin ^2\theta } \frac{\partial ^2}{\partial \phi ^2} 
.\]
This is called the Laplacian on the sphere. For the $\theta ,\phi $ part there is no potential, so it may be more simple than the $r$ part. Our goal is to find eigenfunctions $Y(\theta ,\phi )$ of $\hat{\nabla }^2$, and we can just factor in the $1/r^2$ later. These eigenfunctions are called the \emph{spherical harmonics}.
